% GENERATED FILE. Edit canonical lessons, then run npm run generate:books.
% canonical-source-hash: fnv1a64:8dae8099b528a5ea
% canonical-lessons: BN-C385-ghonta, BN-C385-ghori, BN-C385-tiktik-kora, BN-C385-opekkha-kora, BN-C385-bench

\chapter{Hours and Waiting}
\label{ch:bn-hours-and-waiting}
\hlchaptermodality{385}

\begin{chapteropening}
\textbf{By the end of this chapter:} \emph{I can say how long I waited on a platform bench, and hear the station clock tick.}

Complete the last lesson of chapter 385: I can say how long I waited on a platform bench, and hear the station clock tick.
\end{chapteropening}

\section[\texorpdfstring{\bn{ঘণ্টা} (ghônṭā)}{ঘণ্টা (ghônṭā)}]{\bn{ঘণ্টা} (ghônṭā) — an hour; a bell}

\label{lesson:BN-C385-ghonta}



\begin{quote}
Before the new one: say the Bengali for a local train, then the Bengali for to blow a whistle.
\end{quote}

\subsection*{You'll want to know: \bn{ঘণ্টা}}
\textbf{\bn{ঘণ্টা}} — \emph{ghônṭā} — \textquotedblleft{}an hour; a bell\textquotedblright{}.

\begin{quote}
\bn{ট্রেনটা} \bn{এক} \bn{ঘণ্টা} \bn{দেরিতে} \bn{এল।}
\end{quote}

\begin{quote}
\emph{ṭrenṭā ek ghônṭā derite elô.}
\end{quote}

\begin{quote}
The train came an hour late.
\end{quote}

After a number there is no plural: \bn{দুই} \bn{ঘণ্টা}, two hours.

The first letter is \bn{ঘ}, the breathy partner of \bn{গ}.

\begin{cousinweb}
\textbf{\bn{ঘণ্টা}} (\emph{ghônṭā}) is Sanskrit \emph{ghaṇṭā}, \textquotedblleft{}a bell\textquotedblright{} (Monier-Williams); the sense \textquotedblleft{}an hour\textquotedblright{} is probably the bell or gong that was struck to mark the time.
\end{cousinweb}

\begin{grammarlens}[title={What you've built}]
One more word for getting around.
\end{grammarlens}

\subsection*{Your turn}
\begin{itemize}
\raggedright
  \item \emph{Say it:} \emph{ghônṭā}
  \item \emph{Say it:} \emph{ghônṭā}, once more
  \item \emph{Say it:} \emph{ṭrenṭā ek ghônṭā derite elô.} — The train came an hour late.
  \item \emph{Recall:} say the Bengali for a local train, then the Bengali for to blow a whistle, then say \emph{ghônṭā} again
\end{itemize}

\begin{tcolorbox}[breakable,colback=teal!4,colframe=teal!35!black,title={Before you move on}]
What does \bn{ঘণ্টা} mean? (\textquotedblleft{}An hour; a bell\textquotedblright{}.) Say it once more.
\end{tcolorbox}
\section[\texorpdfstring{\bn{ঘড়ি} (ghôṛi)}{ঘড়ি (ghôṛi)}]{\bn{ঘড়ি} (ghôṛi) — a clock; a watch}

\label{lesson:BN-C385-ghori}



\begin{quote}
Before the new one: say the Bengali for to blow a whistle, then the Bengali for an hour.
\end{quote}

\subsection*{You'll want to know: \bn{ঘড়ি}}
\textbf{\bn{ঘড়ি}} — \emph{ghôṛi} — \textquotedblleft{}a clock; a watch\textquotedblright{}.

\begin{quote}
\bn{আমার} \bn{ঘড়িটা} \bn{হারিয়ে} \bn{গেছে।}
\end{quote}

\begin{quote}
\emph{āmār ghôṛiṭā hāriye gechhe.}
\end{quote}

\begin{quote}
My watch has gone missing.
\end{quote}

One word for the clock on the wall and the watch on your wrist.

\begin{cousinweb}
\textbf{\bn{ঘড়ি}} (\emph{ghôṛi}) goes back to Sanskrit \emph{ghaṭikā}, \textquotedblleft{}a small water-jar\textquotedblright{} (Monier-Williams): a pierced bowl that sank in a set time was an old Indian clock, and the word probably passed from the bowl to the time it kept and then to any clock.
\end{cousinweb}

\begin{grammarlens}[title={What you've built}]
One more word for getting around.
\end{grammarlens}

\subsection*{Your turn}
\begin{itemize}
\raggedright
  \item \emph{Say it:} \emph{ghôṛi}
  \item \emph{Say it:} \emph{ghôṛi}, once more
  \item \emph{Say it:} \emph{āmār ghôṛiṭā hāriye gechhe.} — My watch has gone missing.
  \item \emph{Recall:} say the Bengali for to blow a whistle, then the Bengali for an hour, then say \emph{ghôṛi} again
\end{itemize}

\begin{tcolorbox}[breakable,colback=teal!4,colframe=teal!35!black,title={Before you move on}]
What does \bn{ঘড়ি} mean? (\textquotedblleft{}A clock; a watch\textquotedblright{}.) Say it once more.
\end{tcolorbox}
\section[\texorpdfstring{\bn{টিকটিক} \bn{করা} (ṭikṭik kôrā)}{টিকটিক করা (ṭikṭik kôrā)}]{\bn{টিকটিক} \bn{করা} (ṭikṭik kôrā) — to tick (of a clock)}

\label{lesson:BN-C385-tiktik-kora}



\begin{quote}
Before the new one: say the Bengali for an hour, then the Bengali for a clock.
\end{quote}

\subsection*{You'll want to know: \bn{টিকটিক} \bn{করা}}
\textbf{\bn{টিকটিক} \bn{করা}} — \emph{ṭikṭik kôrā} — \textquotedblleft{}to tick (of a clock)\textquotedblright{}.

\begin{quote}
\bn{স্টেশনের} \bn{বড়} \bn{ঘড়িটা} \bn{টিকটিক} \bn{করছিল।}
\end{quote}

\begin{quote}
\emph{sṭeshôner bôṛo ghôṛiṭā ṭikṭik kôrchhilô.}
\end{quote}

\begin{quote}
The big station clock was ticking.
\end{quote}

Past (he/she …): \bn{সেই} \bn{ঘড়ি} \bn{টিকটিক} \bn{করল}, \emph{sei ghôṛi ṭikṭik kôrlô}.

\bn{টিকটিকি}, a gecko, is named the same way, for its clicking call.

\begin{cousinweb}
\textbf{\bn{টিকটিক} \bn{করা}} (\emph{ṭikṭik kôrā}) joins \bn{টিকটিক}, an imitative word for the tick of a clock, to \bn{করা}, to do. No dictionary we use records an older form, so its history is not settled.
\end{cousinweb}

\begin{grammarlens}[title={What you've built}]
One more word for getting around.
\end{grammarlens}

\subsection*{Your turn}
\begin{itemize}
\raggedright
  \item \emph{Say it:} \emph{ṭikṭik kôrā}
  \item \emph{Say it:} \emph{ṭikṭik kôrā}, once more
  \item \emph{Say it:} \emph{sṭeshôner bôṛo ghôṛiṭā ṭikṭik kôrchhilô.} — The big station clock was ticking.
  \item \emph{Recall:} say the Bengali for an hour, then the Bengali for a clock, then say \emph{ṭikṭik kôrā} again
\end{itemize}

\begin{tcolorbox}[breakable,colback=teal!4,colframe=teal!35!black,title={Before you move on}]
What does \bn{টিকটিক} \bn{করা} mean? (\textquotedblleft{}To tick (of a clock)\textquotedblright{}.) Say it once more.
\end{tcolorbox}
\section[\texorpdfstring{\bn{অপেক্ষা} \bn{করা} (ôpekkhā kôrā)}{অপেক্ষা করা (ôpekkhā kôrā)}]{\bn{অপেক্ষা} \bn{করা} (ôpekkhā kôrā) — to wait}

\label{lesson:BN-C385-opekkha-kora}



\begin{quote}
Before the new one: say the Bengali for a clock, then the Bengali for to tick (of a clock).
\end{quote}

\subsection*{You'll want to know: \bn{অপেক্ষা} \bn{করা}}
\textbf{\bn{অপেক্ষা} \bn{করা}} — \emph{ôpekkhā kôrā} — \textquotedblleft{}to wait\textquotedblright{}.

\begin{quote}
\bn{আমরা} \bn{প্ল্যাটফর্মে} \bn{এক} \bn{ঘণ্টা} \bn{অপেক্ষা} \bn{করলাম।}
\end{quote}

\begin{quote}
\emph{āmrā plæṭphôrme ek ghônṭā ôpekkhā kôrlām.}
\end{quote}

\begin{quote}
We waited on the platform for an hour.
\end{quote}

Past (he/she …): \bn{তিনি} \bn{অপেক্ষা} \bn{করলেন}, \emph{tini ôpekkhā kôrlen}.

The one you wait for takes -\bn{এর} \bn{জন্য}: \bn{ট্রেনের} \bn{জন্য} \bn{অপেক্ষা} \bn{করা}, to wait for the train.

\begin{cousinweb}
\textbf{\bn{অপেক্ষা} \bn{করা}} (\emph{ôpekkhā kôrā}) joins Sanskrit \emph{apekṣā}, \textquotedblleft{}looking round, regard, expectation\textquotedblright{} (Monier-Williams), to \bn{করা}, to do; Bengali narrowed the looking-out to waiting.
\end{cousinweb}

\begin{grammarlens}[title={What you've built}]
One more word for getting around.
\end{grammarlens}

\subsection*{Your turn}
\begin{itemize}
\raggedright
  \item \emph{Say it:} \emph{ôpekkhā kôrā}
  \item \emph{Say it:} \emph{ôpekkhā kôrā}, once more
  \item \emph{Say it:} \emph{āmrā plæṭphôrme ek ghônṭā ôpekkhā kôrlām.} — We waited on the platform for an hour.
  \item \emph{Recall:} say the Bengali for a clock, then the Bengali for to tick (of a clock), then say \emph{ôpekkhā kôrā} again
\end{itemize}

\begin{tcolorbox}[breakable,colback=teal!4,colframe=teal!35!black,title={Before you move on}]
What does \bn{অপেক্ষা} \bn{করা} mean? (\textquotedblleft{}To wait\textquotedblright{}.) Say it once more.
\end{tcolorbox}
\section[\texorpdfstring{\bn{বেঞ্চ} (bench)}{বেঞ্চ (bench)}]{\bn{বেঞ্চ} (bench) — a bench}

\label{lesson:BN-C385-bench}



\begin{quote}
Before the new one: say the Bengali for to tick (of a clock), then the Bengali for to wait.
\end{quote}

\subsection*{You'll want to know: \bn{বেঞ্চ}}
\textbf{\bn{বেঞ্চ}} — \emph{bench} — \textquotedblleft{}a bench\textquotedblright{}.

\begin{quote}
\bn{প্ল্যাটফর্মের} \bn{বেঞ্চে} \bn{বসে} \bn{আমরা} \bn{দুই} \bn{ঘণ্টা} \bn{অপেক্ষা} \bn{করলাম।}
\end{quote}

\begin{quote}
\emph{plæṭphôrmer benche bôse āmrā dui ghônṭā ôpekkhā kôrlām.}
\end{quote}

\begin{quote}
We sat on a bench on the platform and waited for two hours.
\end{quote}

\begin{cousinweb}
\textbf{\bn{বেঞ্চ}} (\emph{bench}) is borrowed from English \emph{bench}; \bn{ঞ্চ} is the n-ch joined letter.
\end{cousinweb}

\begin{grammarlens}[title={What you've built}]
One more word for getting around.
\end{grammarlens}

\subsection*{Your turn}
\begin{itemize}
\raggedright
  \item \emph{Say it:} \emph{bench}
  \item \emph{Say it:} \emph{bench}, once more
  \item \emph{Say it:} \emph{plæṭphôrmer benche bôse āmrā dui ghônṭā ôpekkhā kôrlām.} — We sat on a bench on the platform and waited for two hours.
  \item \emph{Recall:} say the Bengali for to tick (of a clock), then the Bengali for to wait, then say \emph{bench} again
\end{itemize}

\begin{tcolorbox}[breakable,colback=teal!4,colframe=teal!35!black,title={Before you move on}]
What does \bn{বেঞ্চ} mean? (\textquotedblleft{}A bench\textquotedblright{}.) Say it once more.
\end{tcolorbox}
